%! Matrices Modeling Contexts — Unit 4, Lesson 4.14 — Guided Notes Key
\documentclass[10pt]{article}
\usepackage{precalculus-article}
\usepackage{precalculus-key}
\newcommand{\TallMath}[1]{$\displaystyle #1\rule[-1.4em]{0pt}{3.2em}$}

\begin{document}

\pageheader{Unit 4, Lesson 4.14}{Guided Notes Key}
\namedateperiod

\vspace{0.08in}

\begin{objectivebox}
\textbf{4.14.A} \ansline{Build a transition matrix from percent-change data; interpret columns as proportions from each state.}
\textbf{4.14.B} \ansline{Use $T\vec{s}_n=\vec{s}_{n+1}$ to predict future states; find the steady state; use $T^{-1}$ to recover past states.}
\end{objectivebox}

\vspace{0.06in}

\begin{vocabbox}
\termblanklong{transition matrix} \ansline{A square matrix $T$ whose columns give the proportions moving from each state to every other state; columns sum to 1.}
\termblanklong{state vector} \ansline{A column vector $\vec{s}$ recording the count or proportion in each state at a given step.}
\termblanklong{steady state} \ansline{A vector $\vec{s}_\infty$ satisfying $T\vec{s}_\infty=\vec{s}_\infty$; the distribution no longer changes.}
\end{vocabbox}

\vspace{0.06in}

\begin{hookbox}
\textit{A school cafeteria offers two lunch choices every day: \textbf{Pizza (P)} and \textbf{Salad (S)}.
On Day~0, \textbf{100 students} choose Pizza and \textbf{80 students} choose Salad.
Each day, 70\% of Pizza choosers pick Pizza again (30\% switch to Salad), and 60\% of Salad choosers
pick Salad again (40\% switch to Pizza).}

\vspace{6pt}
Prediction --- after many days, approximately how many students will choose each option?

\ansline{Approximately 103 Pizza and 77 Salad (approaches the steady state).}
\end{hookbox}

\vspace{0.06in}

\begin{notesbox}{Building the Transition Matrix}

Each \textbf{column} describes what fraction of members of that state go to each state next.

\vspace{6pt}
\begin{center}
\begin{tabular}{l|cc}
 & \textbf{From P} & \textbf{From S} \\\hline
\textbf{To P} & \ans{0.7} & \ans{0.4} \\[4pt]
\textbf{To S} & \ans{0.3} & \ans{0.6} \\
\end{tabular}
\end{center}

\vspace{8pt}
$T = \begin{bmatrix}\ans{0.7} & \ans{0.4} \\ \ans{0.3} & \ans{0.6}\end{bmatrix}$
\qquad Each column sums to \ans{1}.

\vspace{6pt}
Initial state vector: $\vec{s}_0 = \begin{bmatrix}100 \\ 80\end{bmatrix}$

\end{notesbox}

\vspace{0.06in}

\begin{notesbox}{Predicting Future States}

To find the \textbf{next} state vector, multiply: \quad next state $= \ans{$T$} \times \ans{$\vec{s}_n$}$

\vspace{8pt}
$\vec{s}_1 = T\vec{s}_0
= \begin{bmatrix}0.7(100)+0.4(80) \\ 0.3(100)+0.6(80)\end{bmatrix}
= \ans{$\begin{bmatrix}102\\78\end{bmatrix}$}$

\vspace{10pt}
$\vec{s}_2 = T\vec{s}_1
= \ans{$\begin{bmatrix}0.7(102)+0.4(78)\\0.3(102)+0.6(78)\end{bmatrix}=\begin{bmatrix}102.6\\77.4\end{bmatrix}$}$

\end{notesbox}

\vspace{0.06in}

\begin{notesbox}{The Steady State}

A \textbf{steady state} $\vec{s}_\infty$ satisfies: $T\vec{s}_\infty = \ans{$\vec{s}_\infty$}$

\vspace{6pt}
Repeated multiplication converges: the state vectors approach a fixed vector that never changes.

\vspace{6pt}
What pattern do you notice as you compute more state vectors?
\ansline{The vectors get closer and closer together; the changes shrink; they converge to $\approx\langle102.9,77.1\rangle$.}

\end{notesbox}

\vspace{0.06in}

\begin{notesbox}{Finding Past States}

Since $\vec{s}_1 = T\vec{s}_0$, we can recover $\vec{s}_0$ by multiplying by the \textbf{inverse}:

\vspace{6pt}
$\vec{s}_{n-1} = \ans{$T^{-1}$} \cdot\, \vec{s}_n$

\vspace{8pt}
Apply to find Day~0 from Day~1: $T^{-1}\vec{s}_1 = \ans{$\begin{bmatrix}100\\80\end{bmatrix}$}$

\vspace{6pt}
Interpretation: \ansline{Multiplying by $T^{-1}$ reverses one transition step, recovering the previous state.}

\end{notesbox}

\vspace{0.06in}

\begin{practicebox}
Given $T = \begin{bmatrix}0.6&0.3\\0.4&0.7\end{bmatrix}$ and
$\vec{s}_0 = \begin{bmatrix}50\\50\end{bmatrix}$,
find $\vec{s}_1 = T\vec{s}_0$.

\vspace{8pt}
$\vec{s}_1 = \ans{$\begin{bmatrix}45\\55\end{bmatrix}$}$

\vspace{6pt}
Show work: \ansline{$0.6(50)+0.3(50)=30+15=45$;\quad $0.4(50)+0.7(50)=20+35=55$.}
\end{practicebox}

\end{document}
